%Document Class
\documentclass[a4paper,12pt,reqno]{amsart}


  \headheight0.6in
  \headsep22pt
  \textheight23cm
  \topmargin-1.7cm
  \oddsidemargin 0.5cm
  \evensidemargin0.5cm
  \textwidth15.3cm

%Packages 
\usepackage{amscd,amssymb,amsmath,amsfonts,amsthm, mathrsfs,xspace}
\usepackage{graphicx}
\usepackage{verbatim,enumerate,paralist}
\usepackage{textcomp}
%\usepackage[authoryear]{natbib}
\usepackage[pagebackref,colorlinks]{hyperref}
\hypersetup{%
  urlcolor=blue,linkcolor=blue,citecolor=blue
}
\usepackage{pdfpages}
%\usepackage{pdfsync}%  remove this when finalised

\usepackage[T1]{fontenc}
\usepackage[all,2cell,poly,knot,tips]{xy}
\UseAllTwocells
\CompileMatrices

\def\blue{\color{blue}}
\def\red{\color{red}}
\def\black{\color{black}}

%\usepackage{xypdf} % may not be necessary
\setlength{\marginparwidth}{2cm}
\newcommand{\missing}[1]{\marginpar{\color{red}missing:\\#1}}

%Theoretic Environment Setup
\newcounter{Definitioncount}
\setcounter{Definitioncount}{0}

\newtheorem*{Theorem}{Theorem}% no counter
\newtheorem{theorem}{Theorem}[section] % with counter
\newtheorem{lemma}[theorem]{Lemma}
\newtheorem{proposition}[theorem]{Proposition}
% \newtheorem*{Proposition}{Proposition}
\newtheorem{corollary}[theorem]{Corollary}
% \newtheorem*{Corollary}{Corollary}% no counter

\newtheorem*{Conjecture}{Conjecture}% no counter
\newtheorem*{Lemma}{Lemma}% no counter

\theoremstyle{definition}
\newtheorem*{Definition}{Definition}% no counter
\newtheorem*{Observation}{Observation}% no counter
\newtheorem*{Observations}{Observations}% no counter
\newtheorem*{Remark}{Remark}% no counter
\newtheorem{remark}[theorem]{Remark}
\newtheorem*{Example}{Example}% no counter
\newtheorem{example}[theorem]{Example}

\renewcommand{\theexample}{\arabic{example}}

\newtheoremstyle{fact}{\bigskipamount}{\medskipamount}{\upshape}{}{\itshape}{. }{ }{Fact}
\theoremstyle{fact}
\newtheorem*{Fact}{Fact}% no counter

\newtheoremstyle{genquest}{\bigskipamount}{\medskipamount}{\upshape}{}{\itshape}{. }{ }{General Question}
\theoremstyle{genquest}
\newtheorem*{GenQuest}{General Question}% no counter

% use \newtheoremstyle to underline the header
\newtheoremstyle{step}{2\bigskipamount}{\medskipamount}{\upshape}{}{\itshape}{. }{ }{\underline{Step~\thestep}}
\theoremstyle{step}
\newtheorem{step}{Step}[subsection]
\renewcommand{\thestep}{\arabic{step}}

\newcommand{\lra}{\longrightarrow}
\newcommand{\lla}{\longleftarrow}
\newcommand{\Lra}{\Longrightarrow}
\newcommand{\Lla}{\Longleftarrow}
\newcommand{\ldual}[1]{\mathord{{\let\nolimits\relax\sideset{^\wedge}{}{#1}}}}
\newcommand{\laction}[2]{\mathord{{\let\nolimits\relax\sideset{^{#1}}{}{#2}}}}
\newcommand{\conj}[2]{\mathord{{\let\nolimits\relax\sideset{^{#1}}{}{#2}}}}
\newcommand{\xylongto}[2][2pt]{{\UseTips{}\xymatrix @C+#1{ \ar[r]^{#2}&}}}

%% Script Shortcuts 
\def\CA{{\mathscr A}}
\def\CB{{\mathscr B}}
\def\CC{{\mathscr C}}
\def\CD{{\mathscr D}}
\def\CE{{\mathscr E}}
\def\CF{{\mathscr F}}
\def\CG{{\mathscr G}}
\def\CH{{\mathscr H}}
\def\CI{{\mathscr I}}
\def\CJ{{\mathscr J}}
\def\CK{{\mathscr K}}
\def\CL{{\mathscr L}}
\def\CM{{\mathscr M}}
\def\CN{{\mathscr N}}
\def\CO{{\mathscr O}}
\def\CP{{\mathscr P}}
\def\CQ{{\mathscr Q}}
\def\CR{{\mathscr R}}
\def\CS{{\mathscr S}}
\def\CT{{\mathscr T}}
\def\CU{{\mathscr U}}
\def\CV{{\mathscr V}}
\def\CW{{\mathscr W}}
\def\CX{{\mathscr X}}
\def\CY{{\mathscr Y}}
\def\CZ{{\mathscr Z}}

\DeclareMathOperator{\Lan}{Lan}
%\DeclareMathOperator{\ev}{ev}
\DeclareMathOperator{\Skew}{Skew}
\DeclareMathOperator{\Mnd}{Mnd}
\DeclareMathOperator{\LB}{LBrMon}
\DeclareMathOperator{\Span}{Span}

\newcommand{\K}{\ensuremath{\mathscr K}\xspace}
\newcommand{\M}{\ensuremath{\mathscr M}\xspace}

\renewcommand{\L}{\ensuremath{\LB(\M)}\xspace}

%\newcommand{\Span}{\textnormal{Span}\xspace}

\newcommand{\ox}{\otimes}
\newcommand{\x}{\times}

\newcommand{\op}{\ensuremath{^{\textnormal{op}}}}


\def\theequation{{\thesection.\arabic{equation}}}


\begin{document}

\title{Free skew-monoidal categories, orientals and the Tamari lattice}

\author{Stephen Lack}
\address{Department of Mathematics, Macquarie University, NSW 2109, Australia}
\email{steve.lack@mq.edu.au}

\author{Ross Street}
\address{Department of Mathematics, Macquarie University, NSW 2109, Australia}
\email{ross.street@mq.edu.au}

\thanks{Both authors gratefully acknowledge the support of the Australian Research Council Discovery Grant DP1094883; Lack acknowledges with equal gratitude the support of an Australian Research Council Future Fellowship}

\date{\today}
\maketitle

\section{Introduction}

Monoidal categories \cite{CWM} are important in the large as bases for enriched category theory \cite{KellyBook} 
and other purposes \cite{dlVP}.  
They are also important in the small as mathematical structures themselves: in their abstract form they are pseudomonoids
(or monoidales) in monoidal bicategories \cite{{DaySt1997}, {DMcCS}}. 
It is this small viewpoint which drove the generalization of monoidal category to skew-monoidal category \cite{Szl2012}.
In particular, the relevant mathematical structure there is bialgebroid.
We \cite{{SkMon}, {SkEH}} have developed that viewpoint somewhat in the setting of quantum categories.
However, we also have found \cite{scc} that enriched category theory develops naturally with a skew-monoidal base.

At the outset \cite{{Rice}, {Kelly1964}}, partly as a test for correctness of axioms, lay the question of coherence in monoidal categories. 
That is, which diagrams involving the structural isomorphisms (constraints) commute? 
A deeper understanding \cite{aac} of coherence questions for categorical structures was that they involved knowledge 
of the free structures. 
In some interesting cases (called ``clubable''), it was seen that the free structures on any category could essentially be 
constructed from the free structure on a single generating object (that is, on the one-object discrete category $\mathbf{1}$).
Skew-monoidal categories fall under this classification.

The notion of free needs to be made clear. One meaning is that we consider the underlying 2-functor from the 
2-category of structures, with morphisms preserving structure up to coherent isomorphism, to the 2-category of categories;
the left biadjoint to this gives the free structure on each category. These free structures are only unique up to equivalence.
In the case of monoidal categories, we can take the free structure on $\mathbf{1}$ simply to be the discrete category $\mathbb{N}$
with the natural numbers as objects and addition as the tensor product. One might call this notion of free {\em bicategorical}.   

Another meaning is that we consider the underlying 2-functor from the 2-category of structures, with morphisms {\em strictly}
preserving the structure, to the 2-category of categories; the left adjoint to this gives the free structure on each category. 
These free structures are unique up to isomorphism. This is freeness in the {\em strict} sense. A discussion of these matters can be found in \cite{btc}.

The present paper describes a model for the free skew-monoidal category $\mathrm{Fsk}$ on $\mathbf{1}$ in the strict sense.
Our construction involves the orientals of \cite{aos} and the Tamari lattices of \cite{HT}.
 
\section{Axioms, non-axioms and an important example}

A {\em (left) skew-monoidal category} is a category $\CC$ equipped with an object $I$ (called the {\em unit} or {\em skew unit}),
a functor $\otimes : \CC \times \CC \lra \CC$ (called {\em tensor product}), and
natural families of {\em lax constraints} having the directions 
\begin{equation}
\alpha_{XYZ} : (X\otimes Y)\otimes Z \lra X\otimes (Y\otimes Z)
\end{equation}
\begin{equation}
\lambda_X : I\otimes X \lra X
\end{equation}
\begin{equation}
\rho_X :  X \lra X\otimes I
\end{equation}
subject to five conditions: 

\begin{equation}\label{pentagon}
\xymatrix{
& (W\otimes X)\otimes (Y\otimes Z) \ar[rd]^-{{\alpha_{W,X,Y\otimes Z}}}  & \\
((W\otimes X)\otimes Y)\otimes Z \ar[ru]^-{{\alpha_{W\otimes X,Y,Z}}} \ar[d]_-{{\alpha_{W,X,Y}} \otimes 1_Z} & & W\otimes (X\otimes (Y\otimes Z))  \\
(W\otimes (X\otimes Y))\otimes Z \ar[rr]_-{{\alpha_{W,X\otimes Y,Z}}} & & W\otimes ((X\otimes Y)\otimes Z) \ar[u]_-{1_W\otimes {\alpha_{X, Y,Z}}} }
\end{equation}

\begin{equation}\label{leftunit}
\xymatrix{
(I\otimes X)\otimes Y \ar[rd]_{\lambda_{X} \otimes 1_Y}\ar[rr]^{\alpha_{I,X,Y}}   && I\otimes (X\otimes Y) \ar[ld]^{\lambda_{X\otimes Y}} \\
& X\otimes Y  &
}
\end{equation}

\begin{equation}\label{midunit}
\xymatrix{
(X\otimes I)\otimes Y \ar[rr]^-{\alpha_{X,I,Y}} && X\otimes (I\otimes Y) \ar[d]^-{1_X\otimes \lambda_Y} \\
X\otimes Y \ar[u]^-{\rho_X \otimes 1_Y} \ar[rr]_-{1_{X\otimes Y}} && X\otimes Y}
\end{equation}

\begin{equation}\label{rightunit}
\xymatrix{
(X\otimes Y)\otimes I \ar[rr]^{\alpha_{I,X,Y}}   && X\otimes (Y\otimes I)  \\
& X\otimes Y \ar[lu]^{\rho_{X \otimes Y}}  \ar[ru]_{1_X\otimes \rho_Y}&
}
\end{equation}

\begin{equation}\label{unitunit}
\xymatrix{
I \ar[rd]_{\rho_I}\ar[rr]^{1_I}   && I  \\
& I\otimes I \ar[ru]_{\lambda_I} & .
}
\end{equation}

Notice that we obtain idempotents
$$\epsilon^{\ell}_{X,Y} : (X\otimes I) \otimes Y \stackrel{\alpha_{X,I,Y}} \lra X\otimes (I \otimes Y)
\stackrel{1_X \otimes \lambda_X} \lra X\otimes Y
\stackrel{\rho_X \otimes 1_Y} \lra (X\otimes I)\otimes Y \ ,$$
$$\epsilon^{r}_{X,Y} : X\otimes (I \otimes Y) \stackrel{1_X \otimes \lambda_X} \lra X\otimes Y
\stackrel{\rho_X \otimes 1_Y} \lra (X\otimes I)\otimes Y \stackrel{\alpha_{X,I,Y}} \lra X\otimes (I \otimes Y)$$
and
$$\epsilon_0 : I\otimes I \stackrel{\lambda_I} \lra I \stackrel{\rho_I}\lra I\otimes I $$
which are not identities in general. Moreover, 
$$\alpha_{X,I,Y} : ((X\otimes I) \otimes Y, \epsilon^{\ell}_{X,Y}) \lra (X\otimes (I \otimes Y), \epsilon^{r}_{X,Y})$$
is a morphism of idempotents. 

We shall give an example of a skew-monoidal category that plays an important role in this paper.
Along with it, we fix some notation.
Recall the algebraists' simplicial category $\mathbf{\Delta}$.
The objects are the finite ordinals $$\mathbf{n} = \{0,1,\dots ,n-1\} \ .$$
The morphisms are order-preserving functions $\xi : \mathbf{m} \lra \mathbf{n}$. 
For each such function, we put
$$\xi_{\ell} = \{i\in \mathbf{m} : \xi (i) = \xi (i+1)\}$$
and
$$\xi_{r} = \{j \in \mathbf{n} : j \notin \mathrm{im} \xi \} \ .$$
Notice that $\xi_{\ell} =\varnothing$ means $\xi$ is injective and $\xi_{r} =\varnothing$ means $\xi$ is surjective. 
Then there is a unique factorization $\xi = \partial \circ \sigma$ where 
$\sigma_{\ell} = \xi_{\ell}$, $\sigma_{r} = \varnothing$, $\partial_{\ell} = \varnothing$, and $\partial_{r} = \xi_r$.  

It is well known \cite{osed} that $\mathbf{\Delta}$ is the free strict monoidal category containing a monoid.
The unit object is $\mathbf{0}$ and tensor product is ordinal sum $\mathbf{m}\otimes \mathbf{n} = \mathbf{m+n}$, $\xi\otimes \zeta = \xi \mathbf{+} \zeta$.
The generic monoid is $\mathbf{1}$ with the only multiplication and unit possible.

Let $\mathbf{\Delta}_{\bot}$ denote the subcategory of $\mathbf{\Delta}$ consisting of non-empty finite ordinals $\mathbf{n}$
and first-element-preserving functions; that is, the $\xi$ with $0\notin \xi_r$.
The tensor product on 
$\mathbf{\Delta}$ restricts to $\mathbf{\Delta}_{\bot}$ 
but we replace the unit by the terminal object $\mathbf{1}$.
This gives a skew-monoidal category with identity associativity constraint, where 
$$\lambda_{\mathbf{n}} : \mathbf{1+n} \lra \mathbf{n}$$ is the surjection with 
$(\lambda_{\mathbf{n}})_{\ell} = \{0\}$, and 
$$\rho_{\mathbf{n}} : \mathbf{n} \lra \mathbf{n+1}$$ is the injection with
$(\rho_{\mathbf{n}})_{r} = \{n\}$.

When we write $\mathbf{\Delta}_{\bot}$ in future, we mean it equipped with the above skew-monoidal structure. 

We shall use the convention that each subset $x=\{x_0,x_1, \dots ,x_r\}$ of $\mathbf{n}\in \mathbf{\Delta}$ organizes itself as a list $x=(x_0x_1\dots x_r)$ with $x_0<x_1< \dots <x_r$.  

\section{Motivation, orientals and the Tamari lattice}

Our goal is to discover a nice model for the free skew-monoidal category $\mathrm{Fsk}$ generated by a single object $X$. 
The objects are clearly words in the letters $I$ and $X$, bracketed meaningfully pairwise. 
For example, there is the object 
\begin{equation}
(((X(IX))(XX))X) \ ;
\end{equation}
but we can write this as a triangulated polygon
\begin{equation}\label{hexagon1}
\xymatrix{
& 0 \ar[rr]{} \ar[rrrd]{} \ar[rddd]{}  \ar[ld]_{X}  & & 6  &  \\
1 \ar[rrdd]{} \ar[d]_{I} & & & & 5 \ar[lu]_{X} \\
2  \ar[rrd]_{X} & & & & 4  \ar[u]_{X} \\
& &  3 \ar[rruu]{}  \ar[rru]_{X} &
}
\end{equation}
where the edges $i\lra i+1$ are labelled by the symbols $I$ or $X$.
We can omit the labels $X$ since anything not labelled by $I$ can be understood to be labelled by $X$.

Now we recall the {\em orientals} $\CO_n$ \cite{aos}. This is the free $n$-category on the $n$-simplex.
However, we need no more than the underlying 3-category structure. 
The objects (= 0-cells) of $\CO_n$ are the natural numbers $m$ with $0 \le m \le n$. 
A morphism (=1-cell) $a : p \lra q$ implies $p\le q$ and $a=(a_0 a_1\dots a_r)$ is a subset of $q\smallsetminus p$
with $a_0 =p$; we think of $a$ as the path 
$$p\stackrel{(pa_1)}\lra a_1\stackrel{(a_1a_2)}\lra a_2\lra \dots \stackrel{(a_rq)}\lra q \ .$$
According to \cite{aos} we think of the morphism $a : p \lra q$ as a set of doublets 
$$a=\{(pa_1), (a_1a_2),\dots , (a_rq)\} \ .$$  
For our purposes, we are interested in the morphism $b_n : 0 \lra n$ which is $0\stackrel{(0n)}\lra n$ and
the morphism $e_n : 0 \lra n$ which is the path $0\stackrel{(01)}\lra 1 \stackrel{(12)}\lra 2 \dots  \stackrel{(n-1n)}\lra n$. 

Let $\CT_n = \CO_n (0,n)(b_n,e_n)$ as a 2-category. We shall examine this more explicitly. 

The objects of $\CT_n$ are 2-cells $S:b_n \Lra e_n$ in $\CO_n$.
A description of these, adapted from \cite{aos}, is as follows: 
$S$ is a set of three-element subsets $x=(x_0x_1x_2)$ of $\mathbf{n+1}$ 
satisfying the conditions:
\begin{itemize}
\item[(a)] for $x,y\in S$, if two of the three equations $x_0=y_0$, $x_1=y_1$, $x_2=y_2$ hold then so does the third;
\item[(b)] for $x,y\in S$, at most one of the equations $x_1=y_0$, $x_2=y_1$ holds;
\item[(c)] the function $S\lra \{1,2,\dots ,n-1\}$, taking $x$ to its middle element $x_1$, is a bijection. 
\end{itemize} 
These objects $S$ are in bijection with triangulations of the polygon with $n$ sides.
For example, the triangulation \eqref{hexagon1} of the hexagon has 
\begin{equation}
S=\{(013),(123),(035),(345),(056)\} \ .
\end{equation} 
It will be useful to have some notation involving the inverse $t_S$ of the bijection in (c) above:
\begin{eqnarray}
t_S : \{1,2,\dots ,n-1\} \lra S \nonumber \\
t_S(i) = (\ell_S(i), i , r_S(i)) \ .
\end{eqnarray}
Notice that, in the process, we are defining functions 
\begin{eqnarray}
\ell_S : \{1,2,\dots ,n-1\} \lra \mathbf{n-2}
\end{eqnarray}
and 
\begin{eqnarray}\label{rightfns}
r_S: \{1,2,\dots ,n-1\} \lra \{2,3,\dots ,n\} \ .
\end{eqnarray}
We put 
\begin{equation}
S_{\lambda} = \{i\in \mathbf{n} : \ell_S(i+1)=i \}
\end{equation}
and
\begin{equation}
S_{\rho} = \{i\in \mathbf{n} : r_S(i)=i+1 \} \ .
\end{equation}

Now let us look at the morphisms of $\CT_n$.
As 3-cells in $\CO_n$, they are well defined in \cite{aos} as sets $\theta$ of
four-element subsets $x=(x_0x_1x_2x_3)$ of $\mathbf{n+1}$ satisfying
well formedness and movement conditions; also see \cite{{pc},{pc:c}}.
The indecomposable $\theta : S\lra T$ are precisely the singleton sets $\{ x\}$
such that
\begin{equation}
T = (S\smallsetminus \{(x_0x_1x_2),(x_0x_2x_3)\})\cup \{(x_0x_1x_3),(x_1x_2x_3)\} \ .
\end{equation}
In pictures, this means that $S$ should contain the domain of
$x=(x_0x_1x_2x_3)$ in the diagram \eqref{indec3}
as part of its triangulation and $T$ should be obtained from $S$ by 
replacing that part by the codomain.

\begin{eqnarray}\label{indec3}
\xymatrix{
x_0 \ar[r]^-{} \ar[rd]^-{} \ar[d]_-{} & x_3 \\
x_1 \ar[r]_-{}  & x_2 \ar[u]^-{} }
\qquad
\xymatrix{
\stackrel{x} \lra}
\qquad
\xymatrix{
x_0 \ar[r]^-{} \ar[d]_-{} & x_3 \\
x_1 \ar[r]_-{} \ar[ru]^-{}  & x_2 \ar[u]^-{} }
\end{eqnarray}
In terms of bracketings, it means that $T$ is obtained from $S$ by moving just 
one set of brackets to the right:
$$((W_0((W_1W_2)W_3))W_4) \lra ((W_0(W_1(W_2W_3)))W_4)$$
or
$$(W_0(((W_1W_2)W_3)W_4)) \lra (W_0((W_1(W_2W_3))W_4))$$
where the $W_i$ are meaningfully binarily bracketed words in $I$ and $X$.
So this means that the category $\CT_n$ is the Tamari ordered set 
\cite{{Tamari51},{Tamari62}, {HT}}. We write $S \le T$ for this order.

The following result is immediate on combining \cite{HT} and \cite{aos}.

\begin{proposition}\label{TLatt} For 2-cells $S,T:b_n\Lra e_n : 0\lra n$ in the $n$-category 
$\CO_n$, the following conditions are equivalent:
\begin{itemize}
\item[(i)] $S \le T$;
\item[(ii)] there exists a 3-cell $S\lra T$ in $\CO_n$;
\item[(iii)] for all $x\in S$ and $y\in T$, if $x_1=y_1$ then $x_0\le y_0$; 
\item[(iv)] for all $x\in S$ and $y\in T$, if $x_1=y_1$ then $x_2\le y_2$;
\item[(v)] for all $x\in S$ and $y\in T$, if $x_1=y_1$ then $x_0\le y_0$ and $x_2\le y_2$.
\end{itemize}
The functions $r_S$ for $S:b_n\Lra e_n$ in $\CO_n$ are precisely the functions 
$$f : \{1,2,\dots , n-1\} \rightarrow \{2,3,\dots , n \}$$ 
such that $i < f(i)$, and $ i \le j < f(i)$ implies $f(j) \le f(i)$. 
\end{proposition}

\section{The model}
The objects of the category $\mathrm{Fsk}$ are triples $(m,u,S)$ where $m$
is a positive integer,
$u\subseteq \mathbf{m}$ and $S$ is an object of $\CT_m$.
The idea is that $S$ amounts to a bracketed word in $I$ and $X$ regarded
as a triangulated $m$-sided polygon and $u$ consists of those $i$ for which the 
edge $i\lra i+1$ is labelled by $I$, the other edges $i\lra i+1$ 
being understood to be labelled by $X$. Notice that $m=\#S+1$ 

We will identify two classes of morphisms
\begin{equation}
(m,u,S) \lra (n,v,T)
\end{equation}
in $\mathrm{Fsk}$: those of type 
$\alpha$ and those which are $\alpha$-absent.
The morphisms of {\em type} $\alpha$ have $m=n$, $u=v$ and $S\le T$ 
(see Proposition~\ref{TLatt}); so, given the domain and codomain,
there is at most one.

An {\em $\alpha$-absent} morphism
\begin{equation}
\xi : (m,u,S) \lra (n,v,T)
\end{equation}
in $\mathrm{Fsk}$ is a first-element-and-order-preserving function 
$\xi : \mathbf{m} \lra \mathbf{n}$
such that
\begin{enumerate}
\item $\xi_{\ell} \subseteq u$;
\item $\xi_{r} \subseteq v$;
\item $v=\xi(u\smallsetminus \xi_{\ell})\cup \xi_{r}$; and
\item $T=(\xi \mathbf{+}1_{\mathbf{1}}) (S\smallsetminus t_S \{i+1 : i\in \xi_{\ell}\}) \cup t_T \xi_r$.   
\end{enumerate}
In (4), we need to use the function 
$$\xi \mathbf{+}1_{\mathbf{1}} : \mathbf{m+1} \lra \mathbf{n+1}$$
rather than merely $\xi$ since some of the $x\in S$ could have $x_2 = m$.
  
\begin{lemma}\label{absent}
If $\xi : (m,u,L) \lra (n,v,R)$ and $\zeta : (n,v,T) \lra (k,w,U)$ are $\alpha$-absent
and $R\le T$ then there exist $M$ and $V$
with $L\le M$ and $V\le U$ such that 
$\zeta \circ \xi : (m,u,M) \lra (k,w,V)$ is $\alpha$-absent.
\end{lemma}

In favour of Lemma~\ref{absent}, we have the example 
where $u$ and $w$ are empty, $v=\{1\}$, $L=U=\{(012)\}$, 
$R=\{(012)(023)\}$, $T=\{(013)(123)\}$, $\xi = \partial_1$ and $\zeta =\sigma_1$;
in this case, $L=M$ and $V=U$.
This example concerns axiom \eqref{midunit} for a skew-monoidal category; the other axioms
\eqref{leftunit}, \eqref{rightunit} and \eqref{unitunit} provide examples too.  

A general morphism
\begin{equation}
\xi : (m,u,S) \lra (n,v,T)
\end{equation}
in $\mathrm{Fsk}$ is a first-element-and-order-preserving function
$\xi : \mathbf{m} \lra \mathbf{n}$
such that there are morphisms $(m,u,S)\lra (m,u,L)$ and $(n,v,R)\lra (n,v,T)$
of $\alpha$-type and such that $\xi : (m,u,L) \lra (n,v,R)$ is $\alpha$-absent. 
Composition in $\mathrm{Fsk}$ is composition of functions;
this makes sense in light of Lemma~\ref{absent}. 

There is a faithful functor 
\begin{equation}\label{und}
\mathrm{und} : \mathrm{Fsk} \lra \Delta_{\bot} 
\end{equation}
taking $(m,u,S)$ to $\mathbf{m}$ and $\xi$ to $\xi$. 

Now we describe the skew-monoidal structure on $\mathrm{Fsk}$.
For objects $(m,u,S)$ and $(h,p,U)$, put
\begin{equation}
(m,u,S)\otimes (h,p,U) = (m+h,u\otimes p,S\otimes U) 
\end{equation}
where $$u\otimes p = u\cup (m+v)$$
and $$S\otimes U = S\cup \{(0,m,m+n)\} \cup (m+T) \ .$$
Here $m+v = \{m+i : i\in v\}$ and $m+T = \{(m+i,m+j,m+k) : (i,j\k) \in T\}$.  
For morphisms $\xi : (m,u,S) \lra (n,v,T)$ and $\zeta : (h,p,U) \lra (k,q,V)$ in
$\mathrm{Fsk}$, put
\begin{equation}
\xi \otimes \zeta = \xi \mathbf{+} \zeta : (m+h,u\otimes p,S\otimes U) \lra (n+k,v\otimes q,T\otimes V) \ .
\end{equation}

The associativity constraint
\begin{equation}
\alpha_{(m,u,S),(n,v,T),(k,w,U)} : ((m,u,S)\otimes (n,v,T))\otimes (k,w,U) \rightarrow (m,u,S)\otimes ((n,v,T)\otimes (k,w,U))
\end{equation}
is the 3-cell 
$$\{(0,m,m+n,m+n+k)\} : (S\otimes T)\otimes U \Lra S\otimes (T\otimes U)$$ 
in $\CO_{m+n+k}$. 

The skew unit object is $I = (1,\{0\},\varnothing)$.

The left unit constraint 
\begin{equation}
\lambda_{(m,u,S)} :  (1,\{0\},\varnothing)\otimes (m,u,S) \lra (m,u,S)
\end{equation}
is the surjection called 
$$\lambda_{\mathbf{m}} : \mathbf{1+m} \lra \mathbf{m}$$
in $\Delta_{\bot}$.
 
The right unit constraint 
\begin{equation}
\rho_{(m,u,S)} :  (m,u,S) \lra (m,u,S) \otimes (1,\{0\},\varnothing)
\end{equation}
is the injection called 
$$\rho_{\mathbf{m}} : \mathbf{m} \lra \mathbf{m+1}$$
in $\Delta_{\bot}$.

\begin{theorem}
$\mathrm{Fsk}$ is the free skew-monoidal category on a single generating object.
The unique strict monoidal functor \eqref{und} into $\Delta_{\bot}$, taking the generating object to the skew unit, is faithful. 
\end{theorem}

\section{A bijection}

In order to construct an inverse for a function $\theta$ relating the Tamari lattice to the Catalan monoid, 
we require a construction on certain partial maps.

We will define a new partial map $h^{\prime} : [n] \lra [n-1]$ from a given one $h : [n] \lra [n-1]$.
We will let $x_h \subseteq [n]$ denote the domain of $h$ and use the same symbol $h$ for the function $h: x_h \lra [n-1]$. 
Let $r_h$ be the first $i \in x_h$ with $h(i) > 0$, if it exists.

Obtain $x_{h^{\prime}} \subseteq [n]$ from $x_h$ by deleting precisely $h(r)$ 
of the elements directly preceding $r$ in $x_h$. 
Then define $h^{\prime} : x_{h^{\prime}} \lra [n-1]$ to agree with $h$ on all elements of $x_{h^{\prime}}$
except $r_h$ where $h^{\prime}(r_h) = 0$.

Notice that $r_h < r_{h^{\prime}}$ if the latter exists.   
Put 
\begin{eqnarray*}
\mathbb{T}_n = \{ \text{functions} \  f : \{1,2,\dots , n-1\} \rightarrow \{2,3,\dots , n \} : i < f(i), \ \text{and} \\
  \ i \le j < f(i) \ \text{implies} \ f(j) \le f(i)\} \ .
\end{eqnarray*}

Also put 
$$\mathbb{C}_n = \{ \text{order-preserving functions} \ g : [n-2] \rightarrow [n-2] : g(i) \le i \} \ .$$

There is a function $$\theta : \mathbb{T}_n \lra \mathbb{C}_n $$
defined by $\theta (f) (r) = \#f^{-1} \{2, 3, \dots , r+1\}$ for $0 \le r \le n-2$.

\begin{proposition} $\theta$ is a bijection.
\end{proposition}
\begin{proof}
Without going into the full proof, we define the inverse of $\theta$.
Let $g\in \mathbb{C}_n$.
We recursively define a sequence of functions $h^m : x^m \lra [n-1]$, $1 \le m \le q$, where $q$ is the number of non-empty fibres of $g$. 
Put $x^1 = [n-1]$ and $h^1 = \# g^{-1}(i)$. 
Then $h^{m+1} = (h^m)^{\prime}$ with $x^{m} = x_{h^m}$.
Notice that $$\bigcup_{1\le m \le q} x_m \smallsetminus x_{m+1} = \{1, \dots , n-1\} \ . $$
For $i \in x_m \smallsetminus x_{m+1}$, define $$f(i) = r_{h^m} \ .$$ 
Then $\theta(f) = g$.
\end{proof}

\section{Tamari order} 

The Tamari order on $\mathbb{T}_n$ is simply the pointwise order of functions using the natural order on $\{2,3,\dots , n \}$.  
This obviously transports across the bijection $\theta : \mathbb{T}_n \lra \mathbb{C}_n$ to an order on $\mathbb{C}_n$.
However, this description takes some work to disentangle. On the other hand, it is easy to describe the indecomposable
inequalities $g_1 < g_2$ in $\mathbb{C}_n$.  

\begin{proposition} The Tamari order on $\mathbb{T}_n$ transports across $\theta$ to the partial order on 
$\mathbb{C}_n$ generated by the relation
$g_1 < g_2$ if there exists $0<i<j \le n-1$ such that 
\begin{equation*}
  g_1 (k) =
      \begin{cases}
         g_2 (k) + 1 & \text{if } \  i \le k \le j \ , \\
         g_2 (k) & \text{otherwise }.
       \end{cases}
\end{equation*}
\end{proposition}

Clearly $\mathbb{C}_n$ is a monoid with zero under composition of endofunctions of $[n-2]$.
However, it is not an ordered monoid. To see this, take $g_0$, $g_1$, $g_2$, $g_3$ to be the functions
whose values are $012$, $011$, $002$, $001$, respectively. While $g_0 < g_2$, when we multiply
on the left by $g_1$ we see that $g_1 g_0 = g_1$ and $g_1 g_2 = g_3$; yet $g_1 \nleq g_3$. 

\begin{thebibliography}{00}

\bibitem{DaySt1997} Brian J. Day and Ross Street, \textit{Monoidal bicategories and Hopf algebroids}, Advances in Math. \textbf{129} (1997) 99--157. \label{DaySt1997}

\bibitem{DMcCS} Brian J. Day, Paddy McCrudden and Ross Street, \textit{Dualizations and antipodes}, Applied Categorical Structures \textbf{11} (2003) 229--260. \label{DMcCS]}

\bibitem{HT} Dani\`ele Huguet and Dov Tamari, \textit{Problems of associativity: a simple proof for the lattice property of
systems ordered by a semi-associative law}, J. Comb. Theory (A) \textbf{13} (1972) 7--13.\label{HT}  

\bibitem{btc}  Andr\'e Joyal and Ross Street, \textit{Braided tensor categories}, Advances in Mathematics \textbf{102} (1993) 20--78.\label{btc}

\bibitem{Kelly1964} G. Max Kelly, \textit{On MacLane's conditions for coherence of natural associativities, commutativities, etc.}, Journal of Algebra \textbf{1} (1964) 397--402.\label{Kelly1964}

\bibitem{aac} G. Max Kelly, \textit{An abstract approach to coherence}, Lecture Notes in Math. \textbf{281} (Springer-Verlag, 1972) 106--147.\label{aac}

\bibitem{KellyBook} G. Max Kelly, \textit{Basic concepts of enriched category theory}, London Mathematical Society Lecture Note Series \textbf{64} (Cambridge University Press, Cambridge, 1982). \label{KellyBook}

\bibitem{SkMon} Stephen Lack and Ross Street, \textit{Skew monoidales, skew warpings and quantum categories}, Theory and Applications of Categories \textbf{27(3)} (2012) 27--46.
\label{SkMon}

\bibitem{SkEH} Stephen Lack and Ross Street, \textit{A skew-duoidal Eckmann-Hilton argument and quantum categories}, (submitted; see http://arxiv.org/abs/1210.8192).\label{SkEH}

\bibitem{osed} F. William Lawvere, \textit{Ordinal sums and equational doctrines}, Reprints in Theory and Applications of Categories \textbf{18} (2008) 1--303.\label{osed}

\bibitem{Rice} Saunders Mac Lane, \textit{Natural associativity and commutativity}, Rice Univ. Studies \textbf{49} (1963) 28--46.\label{Rice}

\bibitem{CWM} Saunders Mac Lane, \textit{Categories for the Working Mathematician}, Graduate Texts in Mathematics \textbf{5} (Springer-Verlag, 1971).\label{CWM} 

\bibitem{aos} Ross Street, \textit{The algebra of oriented simplexes}, J. Pure Appl. Algebra \textbf{49} (1987) 283--335.\label{aos}

\bibitem{pc} Ross Street, \textit{Parity complexes}, Cahiers de topologie et g\'eom\'etrie diff\'erentielle cat\'egoriques \textbf{32} (1991) 315--343.\label{pc}

\bibitem{pc:c} Ross Street, \textit{Parity complexes: corrigenda}, Cahiers de topologie et g\'eom\'etrie diff\'erentielle cat\'egoriques \textbf{35} (1994) 359--361.\label{pc:c}

\bibitem{dlVP} Ross Street, \textit{Monoidal categories in, and linking, geometry and algebra}, Bulletin Belgian Math. Society  (to appear; see http://arxiv.org/abs/1201.2991).\label{dlVP}

\bibitem{scc} Ross Street, \textit{Skew-closed categories}, Journal of Pure and Applied Algebra (2012), doi: 10.1016/j.jpaa.2012.09.020; also see http://arxiv.org/abs/1205.6522).\label{scc}

\bibitem{Szl2012} Kornel Szlach\'anyi, \textit{Skew-monoidal categories and bialgebroids}, Advances in Mathematics \textbf{231(3--4)} (2012) 1694--1730.
%arXiv:1201.4981v1 (2012).
\label{Szl2012}

\bibitem{Tamari51} Dov Tamari, \textit{Monoides pr\'eordonn\'es et cha\^ines de Malcev}, 
Doctorat\`es-Sciences Math\'ematiques Th\`ese de Math\'ematiques (Parisa, 1951).\label{Tamari51} 

\bibitem{Tamari62} Dov Tamari, \textit{The algebra of bracketings and their enumeration}, Nieuw Archief voor Wiskunde \textbf{10} (1962) 131--146.\label{Tamari62}




\end{thebibliography}

\end{document}